% !TEX root = ../matan.tex

\section{Признаки сходимости: признак Абеля--Дирихле, примеры}

\subsection*{Преобразование Абеля}

\begin{Lemma}{Преобразование Абеля}{}
	Пусть даны числа $a_{n+1},\dots,a_{n+m}$ и $b_{n+1},\dots,b_{n+m}$. Положим
	$A_n:=0$ и $A_j:=\sum_{k=n+1}^{j} a_k$ для $j=n+1,\dots,n+m$. Тогда
	\[
	\sum_{k=n+1}^{n+m} a_k b_k
	= A_{n+m}\,b_{n+m}
	+\sum_{j=n+1}^{n+m-1} A_j\,(b_j-b_{j+1}).
	\]
\end{Lemma}

\begin{proof}
	Так как $a_j=A_j-A_{j-1}$ (при $j\ge n+1$, с учётом $A_n=0$),
	\[
	\sum_{j=n+1}^{n+m} a_j b_j
	=\sum_{j=n+1}^{n+m} (A_j-A_{j-1}) b_j
	=\sum_{j=n+1}^{n+m} A_j b_j-\sum_{j=n+1}^{n+m} A_{j-1} b_j.
	\]
	Во второй сумме сдвинем индекс $j\mapsto j+1$:
	$\sum_{j=n+1}^{n+m} A_{j-1} b_j=\sum_{j=n}^{n+m-1} A_j b_{j+1}
	=\sum_{j=n+1}^{n+m-1} A_j b_{j+1}$ (слагаемое с $j=n$ нулевое, так как
	$A_n=0$). Поэтому
	\[
	\sum_{j=n+1}^{n+m} a_j b_j
	=A_{n+m} b_{n+m}+\sum_{j=n+1}^{n+m-1} A_j b_j-\sum_{j=n+1}^{n+m-1} A_j b_{j+1}
	=A_{n+m} b_{n+m}+\sum_{j=n+1}^{n+m-1} A_j (b_j-b_{j+1}). \qedhere
	\]
\end{proof}

\begin{Remark}{}{}
	Если $(b_k)$ монотонна, то все разности $b_j-b_{j+1}$ имеют один знак,
	поэтому
	$\sum_{j=n+1}^{n+m-1}|b_j-b_{j+1}|=\bigl|\sum (b_j-b_{j+1})\bigr|
	=|b_{n+1}-b_{n+m}|\le|b_{n+1}|+|b_{n+m}|$. Это «телескопирование» полного
	изменения --- ключ к обеим оценкам ниже.
\end{Remark}

\subsection*{Признак Дирихле}

\begin{Theorem}{Признак Дирихле}{}
	Пусть частичные суммы ряда $\sum_{k=1}^{\infty} a_k$ ограничены: существует
	$C_0$ такое, что $\bigl|\sum_{k=1}^{N} a_k\bigr|\le C_0$ для всех $N$. Пусть
	последовательность $(b_k)$ монотонна и $b_k\to0$. Тогда ряд
	$\sum_{k=1}^{\infty} a_k b_k$ сходится.
\end{Theorem}

\begin{proof}
	Из ограниченности частичных сумм следует ограниченность всех «хвостовых»
	сумм: для $j>n$ имеем $A_j=\sum_{k=n+1}^{j} a_k=S_j-S_n$, поэтому
	$|A_j|\le|S_j|+|S_n|\le 2C_0=:C$.

	Зафиксируем $\eps>0$. Так как $b_k\to0$, найдётся $n_0$ такое, что при
	$k\ge n_0$ выполнено $|b_k|<\eps$. Возьмём $n\ge n_0$ и применим
	преобразование Абеля к хвосту:
	\[
	\left|\sum_{k=n+1}^{n+m} a_k b_k\right|
	\le |A_{n+m}|\,|b_{n+m}|
	+\sum_{j=n+1}^{n+m-1} |A_j|\,|b_j-b_{j+1}|.
	\]
	Оценивая $|A_j|\le C$ и используя монотонность $(b_k)$
	($\sum|b_j-b_{j+1}|=|b_{n+1}-b_{n+m}|\le|b_{n+1}|+|b_{n+m}|<2\eps$), получаем
	\[
	\left|\sum_{k=n+1}^{n+m} a_k b_k\right|
	\le C\,|b_{n+m}|+C\bigl(|b_{n+1}|+|b_{n+m}|\bigr)
	< C\eps+C\cdot 2\eps=3C\eps.
	\]
	Так как это верно при всех $m$ и $\eps$ произвольно, по критерию Коши ряд
	$\sum a_k b_k$ сходится.
\end{proof}

\subsection*{Признак Абеля}

\begin{Theorem}{Признак Абеля}{}
	Пусть ряд $\sum_{k=1}^{\infty} a_k$ сходится, а последовательность $(b_k)$
	монотонна и ограничена. Тогда ряд $\sum_{k=1}^{\infty} a_k b_k$ сходится.
\end{Theorem}

\begin{proof}
	Пусть $|b_k|\le B$ для всех $k$. Так как $\sum a_k$ сходится, по критерию
	Коши для любого $\eps>0$ найдётся $n_0$ такое, что при $n\ge n_0$ все
	хвостовые суммы $A_j=\sum_{k=n+1}^{j} a_k$ удовлетворяют $|A_j|<\eps$.
	Применяя преобразование Абеля и монотонность $(b_k)$, имеем при $n\ge n_0$
	\[
	\left|\sum_{k=n+1}^{n+m} a_k b_k\right|
	\le |A_{n+m}|\,|b_{n+m}|+\sum_{j=n+1}^{n+m-1}|A_j|\,|b_j-b_{j+1}|
	\le \eps B+\eps\bigl(|b_{n+1}|+|b_{n+m}|\bigr)
	\le \eps B+\eps\cdot 2B=3B\eps.
	\]
	По критерию Коши ряд $\sum a_k b_k$ сходится.
\end{proof}

\begin{Example}{Применение признака Дирихле}{}
	Ряд $\displaystyle\sum_{k=1}^{\infty}\frac{\sin k}{k}$ сходится. Здесь
	$a_k=\sin k$, $b_k=\frac1k$. Последовательность $b_k=\frac1k$ монотонно
	убывает к нулю, а частичные суммы $\sum_{k=1}^{N}\sin k$ ограничены: по
	формуле суммирования синусов
	\[
	\sum_{k=1}^{N}\sin k=\frac{\sin\frac{N}{2}\,\sin\frac{N+1}{2}}{\sin\frac12},
	\qquad\text{откуда}\qquad
	\left|\sum_{k=1}^{N}\sin k\right|\le\frac{1}{\sin\frac12}.
	\]
	По признаку Дирихле ряд сходится. Аналогично сходится
	$\sum_{k=1}^{\infty}\frac{\cos k}{k}$.
\end{Example}

\begin{Example}{Применение признака Абеля}{}
	Если ряд $\sum a_k$ сходится, то сходится и ряд
	$\sum a_k\bigl(1+\tfrac1k\bigr)^{-1}$, поскольку
	$b_k=\bigl(1+\tfrac1k\bigr)^{-1}$ монотонно возрастает и ограничена
	(стремится к $1$). Здесь сходимость самого $\sum a_k$ существенна --- именно
	она требуется в признаке Абеля.
\end{Example}
